\section{设计技术指标}

本次设计任务要求的指标如下。指标以进一步优化时生成的Zemax优化向导文件为主要依据，并结合最终系统的孔径、视场和材料设置进行整理：

\begin{longtable}[]{@{}
  >{\raggedright\arraybackslash}p{(\columnwidth - 4\tabcolsep) * \real{0.3545}}
  >{\raggedright\arraybackslash}p{(\columnwidth - 4\tabcolsep) * \real{0.3117}}
  >{\raggedright\arraybackslash}p{(\columnwidth - 4\tabcolsep) * \real{0.3338}}@{}}
\toprule\noalign{}
\endhead
\bottomrule\noalign{}
\endlastfoot
\textbf{参数名称} & \textbf{设计指标} & \textbf{说明} \\
有效焦距（EFL） & 30 mm & 优化向导中EFFL目标值为30 mm，实际值约为30.000004 mm \\
最大半视场角 & \(\pm 22.5^\circ\) & 采用5个归一化视场进行像质均衡 \\
入瞳直径 & 4 mm & 光阑半口径为2 mm，与目视系统出瞳尺寸匹配 \\
像质评价指标 & RMS spot x+y centroid & 优化向导采用质心参考的RMS点列半径评价 \\
视场采样 & 0、0.25、0.50、0.75、1.00 & 对应轴上到全视场的5个归一化视场 \\
工作波长 & 486.1 / 587.6 / 656.3 nm & F、d、C三线，主波长为587.6 nm \\
系统总长控制 & 约46.2 mm & 优化向导中TOTR记录值为46.197 mm，用于控制结构紧凑性 \\
光程相关约束 & OPLT=60，OPGT=7 & 优化向导中作为辅助约束，权重均为1 \\
厚度边界约束 & 空气间隔\(\ge\)0.3 mm，玻璃中心/边缘厚度\(\ge\)1.0 mm & 玻璃中心厚度上限为6.0 mm，保证加工可行性 \\
镜片数量 & 3片 & 1片场镜+2片胶合组结构 \\
玻璃材料 & REALVIEW\_1.9\_\newline LIGHTWEIGHT、\newline N-BASF64、H-ZLAF75C & 按进一步优化后的Lens Data设置 \\
\end{longtable}

注意事项：

\begin{enumerate}
\def\labelenumi{\arabic{enumi}.}
\item
  目视光学系统（望远镜目镜）的入瞳即物镜的出瞳，大小通常约为2～5 mm，本设计取4 mm；
\item
  有效焦距按30 mm控制，优化函数中使用EFFL操作数作为主要焦距约束；
\item
  三色波长中以587.6 nm（氦d线）为主波长，对F光（486.1 nm）和C光（656.3 nm）进行色差校正；
\item
  视场设置采用0、0.25、0.50、0.75、1.00五个归一化视场，比仅取轴上、中间和边缘视场更有利于控制全视场像质；
\item
  优化过程中以RMS spot x+y centroid作为主要评价函数，并通过空气、玻璃厚度边界约束防止出现不可加工结构。
\end{enumerate}
